% ---
% id: transmitter/ass:ens
% kind: ass
% title: "Emission of the specified ensemble"
% summary: >-
%   The physical source draws each round's letter independently with the
%   declared weights and emits the product of the declared states, no round
%   depending on another. It ties a source specification declared as matrices
%   to the device it describes.
% prominence: standard
% status: accepted
% uses: [def:source-spec]
% assumes: []
% refs:
%   - "tupkary2026decoy, Protocol 2, step 1(a)"
% provenance: newly-derived
% short_id: A-ENS
% verification: never
% category: source
% presupposes: [trusted-source]
% affects: [secrecy]
% user_statement: >-
%   Your source prepares, in every round, one of the declared states with the
%   declared probabilities. Each round's letter is drawn independently with
%   weights p_x, and the states emitted over the n rounds form a product: no
%   round's letter or state depends on another round, its detection results or
%   its public messages. A source with memory or drift between rounds (for
%   example pattern effects in a modulator) violates this. Where your source is
%   a pair source inside your laboratory with one half measured there to herald
%   or fix the letter, those local detectors belong to the declared ensemble:
%   their efficiency mismatch must be carried by the declared weights and
%   states, and their memory across rounds (for example afterpulsing or dead
%   time) violates this like any other dependence between rounds.
% lean:
% notes: "A specification declared by a source model carries the physical-source assumption its submodule declares in the place of this one (def:source-spec (b); decoy/ass:wcp)."
% ---

The assumption is about the source as it runs: in every round it emits
exactly the declared ensemble, not an approximation of it, and it does so
afresh, independently of every other round. The published prepare-and-measure
protocol makes the same premise, Alice preparing each round's state
``independently for each round'' \cite[Protocol~2, step~1(a)]{tupkary2026decoy};
the product it gives the $n$ rounds is the state from which the
source-replacement lemma starts. A specification declared by a \tref{source
model} carries the premise its submodule declares in the place of this one. In
\tref{untrusted-source mode} the protocol compiler's assumption on Alice's
measurement governs the measurement by which Alice obtains a letter from a
received system.

\begin{assumption}
\label{ass:ens}
For the declared signal ensemble $\{p_x,\rho_x\}$ of
Definition~\ref{def:source-spec} on the \tref{transmitted system}, the
letters the physical source draws and the systems it emits satisfy the
following; a letter that Alice obtains in untrusted-source mode by measuring a
received system is not drawn by the source, and nothing below concerns it:
(i) the letters $x_1,\dots,x_n$ it draws in the $n$ rounds are drawn
independently, each with the weights $p_x$;
(ii) given those letters, the systems it emits in the $n$ rounds are in the
state $\rho_{x_1}\otimes\dots\otimes\rho_{x_n}$;
(iii) neither the letter it draws in a round nor the state it emits in it
depends on what happens in the other rounds, their outcomes and public
messages included.
\end{assumption}
